À l'instant $t_{0}=0$, on prépare une solution ($S$) en mélangeant un volume
d'une solution aqueuse d'iodure de potassium
($\ce{K^+_{(aq)} + I^-_{(aq)}}$)
contenant $n_{1}=\qty{8e-2}{\mol}$ d'ions $\ce{I^-_{(aq)}}$, avec
un volume d'une solution aqueuse de peroxodisulfate de sodium
($\ce{2Na^+_{(aq)} + S2O^2-_{8(aq)}}$)
contenant $n_{2}=\qty{2e-2}{\mol}$ d'ions $\ce{S2O^2-_{8(aq)}}$. Le volume
totale de la solution est $V=\qty{200}{\mL}$. Au cours de
la réaction, il se forme le diiode selon l'équation-bilan~:
\[
    \ce{S2O^2-_{8(aq)} + 2I^-_{(aq)} -> 2SO^2-_{4(aq)} + I2_{(aq)}}
\]

\begin{enumerate}
    \item Déterminer la valeur de l'avancement maximale $x_{\max}$. Déduire le
          réactif limitant.
    \item La courbe de la figure ci-dessous donne les variations de la quantité
          de matière du diiode formé en fonction du temps
          $n(\ce{I2})=f(t)$.
          \begin{figure}[H]
              \centering
              \includegraphics[width=0.7\textwidth]{2025_11_26_673124dc142c0df571a8g-2}
          \end{figure}
          
          \begin{enumerate}
              \item Calculer en (\unit{\mol\per\L\per\min}), la valeur de la
                    vitesse volumique de réaction à l'instant $t_{0}=0$.
              \item La valeur de la vitesse volumique de la réaction à l'instant
                    ${t_{1}=\qty{18}{\min}}$ est
                    ${v_{1}=\qty{1.44e-3}{\mol\per\L\per\min}}$.

                    Expliquer la diminution de la vitesse volumique de la
                    réaction.
              \item Citer un facteur cinétique permettant d'augmenter la vitesse
                    volumique de réaction sans changer l'état initial du système
                    chimique.
              \item Déterminer graphiquement le temps de demi-réaction $t_{1/2}$.
          \end{enumerate}
\end{enumerate}
